% =====================================================================
%  Lecture 23 -- User Plane                               (Persian edition)
% =====================================================================
\documentclass[11pt,a4paper]{article}
\providecommand\ACNroot{../..}
\input{\ACNroot/templates/lecture-fa}

\ACNset{number}{23}
\ACNset{title}{صفحهٔ کاربر نسل پنجم}
\ACNset{shorttitle}{صفحهٔ کاربر}
\ACNset{titleother}{The 5G User Plane}
\ACNset{part}{بخش پنجم -- صفحه‌ها و چشم‌انداز (درس‌های ۲۳ تا ۲۶)}
\ACNset{time}{۶ تا ۸ ساعت (۳ ساعت مطالعه، ۳ تا ۵ ساعت تمرین)}
\ACNset{level}{پیشرفته}
\ACNset{prereq}{درس‌های ۶، ۱۳، ۱۶، ۱۹ و ۲۲}
\ACNset{gen}{5}

\begin{document}
\ACNcover

% ---------------------------------------------------------------------
\section{اهداف یادگیری}
پس از مطالعهٔ این درس باید بتوانید:
\begin{objectives}
  \item پشتهٔ پروتکل صفحهٔ کاربر نسل پنجم را از \acro{UE} تا شبکهٔ داده، همراه با تقسیم \acro{CU}/\acro{DU} در \acro{gNB}، \textbf{رسم کنید}.
  \item نقش \acro{SDAP}، \acro{PDCP}، \acro{RLC} و \acro{MAC} و نیز \acro{GTP-U} روی \lr{F1-U}، \lr{N3} و \lr{N9} را \textbf{توضیح دهید}.
  \item تقسیم‌های کارکردی \acro{RAN} را \textbf{مقایسه کنید} و نیاز فرانت‌هال آن‌ها را \textbf{برآورد کنید}.
  \item پردازش بسته در \acro{UPF} با قواعد \acro{PFCP} و گزینه‌های پیاده‌سازی آن را \textbf{شرح دهید}.
  \item سربار سرآیند، بودجهٔ تأخیر، قابلیت اطمینان تکرار و پنجرهٔ شمارهٔ ترتیب را \textbf{محاسبه کنید}.
\end{objectives}

% ---------------------------------------------------------------------
\section{دانش پیش‌نیاز}
به پشتهٔ پروتکل رادیویی \acro{LTE} (درس ۶)، \acro{GTP-U} و \acro{PFCP} (درس ۱۳)، فنون صفحهٔ دادهٔ \acro{NFV} (درس ۱۶)، توپولوژی‌های
\acro{UPF} (درس ۱۹) و مدل \acro{QoS} (درس ۲۲) نیاز دارید. این درس و درس بعد سامانهٔ نسل پنجم را «بر حسب صفحه» بررسی می‌کنند: نخست مسیر
دادهٔ کاربر، سپس سیگنالینگی که آن را کنترل می‌کند.

\paragraph{واژه‌های کلیدی} واژگان این درس همراه با معادل انگلیسی آن‌ها:
\begin{keyterms}
صفحهٔ کاربر & \lr{User plane} & کارکردها و پروتکل‌هایی که دادهٔ کاربر را حمل می‌کنند. \\
واحد مرکزی / توزیع‌شده / رادیویی & \lr{CU / DU / RU} & بخش‌های یک \acro{gNB} تقسیم‌شده. \\
واسط‌های صفحهٔ کاربر & \lr{F1-U / N3 / N9} & واسط‌های \acro{GTP-U} میان \acro{CU} و \acro{DU}، \acro{RAN} و \acro{UPF}، و دو \acro{UPF}. \\
تقسیم کارکردی & \lr{Functional split} & نقطه‌ای از پشتهٔ پروتکل که \acro{RAN} در آن تقسیم می‌شود. \\
فرانت‌هال & \lr{Fronthaul} & انتقال میان \acro{DU} و \acro{RU} (یا باندپایه و رادیو). \\
تکرار \acro{PDCP} & \lr{PDCP duplication} & ارسال یک \acro{PDU} \acro{PDCP} روی دو مسیر برای قابلیت اطمینان. \\
پنجرهٔ شمارهٔ ترتیب & \lr{SN window} & بازهٔ شماره‌های ترتیبِ در جریان برای بازترتیب‌دهی. \\
قواعد \acro{PFCP} & \lr{PDR / FAR / QER / URR} & قواعد تشخیص، هدایت، \acro{QoS} و مصرف در \acro{UPF}. \\
شتاب صفحهٔ داده & \lr{DPDK / SmartNIC} & شتاب‌دهی نرم‌افزاری و سخت‌افزاری پردازش بسته. \\
خروج محلی & \lr{Local breakout} & خروج به یک شبکهٔ دادهٔ نزدیک به \acro{RAN}. \\
\end{keyterms}

% ---------------------------------------------------------------------
\section{مقدمه و انگیزه}
همهٔ اجزای دیگر یک شبکهٔ موبایل در خدمت صفحهٔ کاربرند: مسیری که بسته‌های مشترک را حمل می‌کند. کارایی آن -- گذردهی، تأخیر، قابلیت اطمینان،
هزینه به ازای هر بیت -- همان چیزی است که کاربران تجربه می‌کنند و اپراتورها برایش هزینه می‌پردازند. در نسل پنجم صفحهٔ کاربر از همیشه
انعطاف‌پذیرتر است: \acro{gNB} را می‌توان به واحدهای متمرکز و توزیع‌شده تقسیم کرد، صفحهٔ کاربر هسته یک \acro{UPF} برنامه‌پذیر است که
می‌تواند در هر جایی از یک مرکز دادهٔ مرکزی تا یک کارخانه قرار گیرد، و چند \acro{UPF} می‌توانند در یک نشست زنجیر شوند.

این انعطاف پرسش‌های طراحی‌ای پدید می‌آورد که این درس شما را برای پاسخ کمّی به آن‌ها آماده می‌کند: با فرانت‌هال موجود \acro{RAN} را کجا تقسیم
کنیم، با یک بودجهٔ تأخیر \acro{UPF}ها را کجا قرار دهیم، چگونه با تکرار به قابلیت اطمینان \acro{URLLC} برسیم، و سرآیندهای پروتکل چه هزینه‌ای
دارند \cite{ts38300,ts38401,ts23501}.

% ---------------------------------------------------------------------
\section{مفاهیم بنیادی}

\subsection{پشتهٔ سرتاسری}
دادهٔ کاربر این مسیر را می‌پیماید (\cref{fig:stack}): در \acro{UE} از \acro{SDAP}، \acro{PDCP}، \acro{RLC}، \acro{MAC} و \acro{PHY}؛ از راه رادیو
به \lr{gNB-DU} (\acro{RLC}، \acro{MAC}، \acro{PHY})؛ روی \lr{F1-U} در \acro{GTP-U} به \lr{gNB-CU-UP} (\acro{PDCP}، \acro{SDAP})؛ روی \lr{N3} در
\acro{GTP-U} به \acro{UPF} و احتمالاً روی \lr{N9} به \acro{UPF}های دیگر؛ و سرانجام از \lr{N6} به شبکهٔ داده. لایه‌های رادیویی این نقش‌ها را
دارند \cite{ts38300}:
\begin{itemize}
  \item \textbf{\acro{SDAP}}: جریان‌های \acro{QoS} را به حامل‌های دادهٔ رادیویی نگاشت و \acro{QFI}ها را علامت‌گذاری می‌کند (درس ۲۲).
  \item \textbf{\acro{PDCP}}: فشرده‌سازی سرآیند (\acro{ROHC})، رمزنگاری و حفاظت یکپارچگی اختیاری، بازترتیب‌دهی و حذف تکراری‌ها، و تکرار
        \acro{PDCP}.
  \item \textbf{\acro{RLC}}: قطعه‌بندی و، در حالت تصدیق‌شده، ارسال مجدد \acro{ARQ}؛ \acro{RLC} در \acro{NR} دیگر الحاق یا بازترتیب‌دهی نمی‌کند.
  \item \textbf{\acro{MAC}}: تسهیم کانال‌های منطقی، \acro{HARQ} و زمان‌بندی.
\end{itemize}

\begin{acnfigure}
  \resizebox{\linewidth}{!}{\acnfig{lecture-23/up-stack}}
  \caption{پشتهٔ پروتکل صفحهٔ کاربر نسل پنجم با یک \acro{gNB} تقسیم‌شده (\lr{CU-UP} و \acro{DU}).}
  \label{fig:stack}
\end{acnfigure}

\subsection{تقسیم \lr{gNB}}
\lr{3GPP} یک \emph{تقسیم لایهٔ بالاتر} را میان یک واحد مرکزی (\acro{CU}: \acro{RRC}، \acro{SDAP}، \acro{PDCP}) و یک واحد توزیع‌شده (\acro{DU}:
\acro{RLC}، \acro{MAC}، \acro{PHY}) استاندارد کرده که با \lr{F1} به هم متصل‌اند \cite{ts38401}؛ \acro{CU} خود به صفحهٔ کنترل (\lr{CU-CP}) و صفحهٔ
کاربر (\lr{CU-UP}) تقسیم می‌شود که با \lr{E1} به هم متصل‌اند. تقسیم‌های لایهٔ پایین‌تر میان \acro{DU} و یک واحد رادیویی (\acro{RU}) -- به‌ویژه
گزینهٔ ۷٫۲ که اتحادیهٔ \lr{O-RAN} مشخص کرده -- بخشی از لایهٔ فیزیکی را به \acro{RU} منتقل می‌کنند (\cref{fig:splits}). هرچه تقسیم پایین‌تر
باشد، کارکردهای بیشتری متمرکز می‌شوند (سایت‌های ساده‌تر، هماهنگی میان سلول‌ها)، اما نرخ فرانت‌هال بالاتر و نیاز تأخیر سخت‌گیرانه‌تر
می‌شود (\cref{sec:fh}).

\begin{acnfigure}
  \resizebox{\linewidth}{!}{\acnfig{lecture-23/splits}}
  \caption{برخی تقسیم‌های کارکردی \acro{RAN} با نرخ‌های تقریبی فرانت‌هال (فرض‌ها در \cref{sec:fh}).}
  \label{fig:splits}
\end{acnfigure}

\subsection{\lr{UPF}}
\acro{UPF} هر بسته را بر اساس قواعدی که \acro{SMF} نصب کرده پردازش می‌کند (\cref{fig:upf}) \cite{ts29244}: یک \emph{\acro{PDR}} بسته را تشخیص
می‌دهد (با \acro{TEID}، نشانی \acro{UE}، فیلتر \acro{SDF} یا \acro{QFI}؛ از میان \acro{PDR}های منطبق، آن‌که بالاترین تقدم را دارد برنده است) و به یک
\emph{\acro{FAR}} (هدایت با سرآیند بیرونی جدید، بافرکردن، دورریختن، تکرار)، \emph{\acro{QER}}ها (دروازه‌بانی و محدودیت نرخ) و
\emph{\acro{URR}}ها (شمارش و گزارش حجم، زمان و رویدادها) اشاره می‌کند. بافرکردن برای \acro{UE}های بیکار از یک \emph{\acro{BAR}} پیروی می‌کند.
بدین‌سان \acro{UPF} کارکردهای صفحهٔ کاربر \lr{S-GW} و \lr{P-GW} در \acro{EPC} را (درس ۱۳) یکجا دارد.

\begin{acnfigure}
  \resizebox{\linewidth}{!}{\acnfig{lecture-23/upf-pipeline}}
  \caption{پردازش بسته در \acro{UPF} با قواعد \acro{PFCP}.}
  \label{fig:upf}
\end{acnfigure}

% ---------------------------------------------------------------------
\section{توضیحات فنی تفصیلی}

\subsection{تکرار \lr{PDCP} و چنداتصالی}
برای \acro{URLLC}، \acro{PDCP} می‌تواند بسته‌ها را تکرار کند و نسخه‌ها را روی دو مسیر مستقل بفرستد: دو حامل از یک \acro{gNB} (تجمیع حامل) یا دو
\acro{gNB} (اتصال دوگانه). \acro{PDCP} گیرنده نسخهٔ دوم را دور می‌ریزد. تکرار، استقلال مسیرها را به قابلیت اطمینان تبدیل می‌کند (\cref{sec:dup})
و وقفهٔ دست‌به‌دست را حذف می‌کند، به بهای دو برابر شدن منابع رادیویی مصرفی برای جریان تکرارشده. این قابلیت را می‌توان هنگام بدترشدن کانال
با عناصر کنترلی \acro{MAC} به‌صورت پویا فعال کرد.

\subsection{شمارهٔ ترتیب و بازترتیب‌دهی}
\acro{PDCP} هر \acro{SDU} را شماره‌گذاری می‌کند (\acro{SN} ۱۲ یا ۱۸ بیتی)؛ \acro{RLC} در حالت تصدیق‌شده قطعه‌ها را شماره‌گذاری می‌کند (\acro{SN} ۱۲ یا ۱۸
بیتی) و در حالت تصدیق‌نشده ۶ یا ۱۲ بیتی. \acro{PDCP} گیرنده به‌ترتیب تحویل می‌دهد و تا پایان یک زمان‌سنج بازترتیب‌دهی منتظر بسته‌های گمشده
می‌ماند. پنجرهٔ \acro{SN}هایی که می‌توانند در جریان باشند نصف فضای \acro{SN} است؛ در نرخ‌های بالا یک \acro{SN} کوتاه به‌سرعت دور می‌زند، و به
همین دلیل \acro{NR} \acro{SN}های ۱۸ بیتی را معرفی کرد (\cref{sec:sn}).

\subsection{انتقال \lr{F1-U} و \lr{N3}}
روی \lr{F1-U}، \lr{CU-UP} و \acro{DU} \acro{PDU}های \acro{PDCP} را در \acro{GTP-U} همراه با سرآیند پروتکل صفحهٔ کاربر \acro{NR} مبادله می‌کنند که
اطلاعات کنترل جریان (مثلاً اندازهٔ بافر مطلوب \acro{DU}) را حمل می‌کند تا \acro{CU} بتواند آهنگ \acro{DU} را تنظیم کند و هنگام تغییر شرایط
رادیویی از سرریز جلوگیری کند. روی \lr{N3} و \lr{N9} سرآیند \acro{GTP-U} در یک سرآیند الحاقی \acro{QFI} و نشانگر \acro{QoS} بازتابی را حمل
می‌کند. هر دو روی شبکه‌های انتقال \acro{IP} اجرا می‌شوند که باید با نگاشت \acro{QFI}ها/\acro{5QI}ها به مقادیر \acro{DSCP} کیفیت خدمات را حفظ کنند.

\subsection{پیاده‌سازی و جانمایی \lr{UPF}}
یک \acro{UPF} باید هر بسته را با نرخ خط طبقه‌بندی، شمارش و دوباره کپسوله کند. پیاده‌سازی‌ها از نرم‌افزار روی کارسازها با \acro{DPDK} و
\lr{SR-IOV} (درس ۱۶)، تا واگذاری کپسوله‌سازی و طبقه‌بندی به \lr{SmartNIC}، تا سوئیچ‌های برنامه‌پذیری که خط لولهٔ \lr{P4} را برای مسیر سریع
اجرا می‌کنند (درس ۱۵) و پردازنده استثناها را رسیدگی می‌کند، گسترده‌اند. جانمایی از بودجهٔ تأخیر پیروی می‌کند (\cref{sec:budget}): \acro{UPF}های
مرکزی برای ترافیک اینترنت، \acro{UPF}های منطقه‌ای یا لبه برای ترافیک حساس به تأخیر و ترافیکی که محلی خدمت‌رسانی می‌شود، و \acro{UPF}های
درون‌سازمانی برای شبکه‌های سازمانی و صنعتی.

% ---------------------------------------------------------------------
\section{مدل‌های ریاضی}

\subsection{سربار پروتکل}\label{sec:oh}
\paragraph{مدل} یک بستهٔ \acro{IP} با $L$ بایت روی رادیو سرآیندهای \acro{SDAP}، \acro{PDCP}، \acro{RLC} و \acro{MAC} به اندازهٔ
$H = H_{\mathrm{SDAP}} + H_{\mathrm{PDCP}} + H_{\mathrm{RLC}} + H_{\mathrm{MAC}}$ بایت می‌گیرد. بازدهی برابر است با
\begin{equation}\label{eq:overhead}
  \eta = \frac{L}{L + H} .
\end{equation}
\paragraph{مثال} \acro{SDAP} ۱، \acro{PDCP} ۳ (\acro{SN} ۱۸ بیتی)، \lr{RLC AM} ۳ و زیرسرآیند \acro{MAC} ۳ بایت: $H = 10$ بایت. برای $L = 1500$:
$\eta = 99.3\,\%$؛ برای یک بستهٔ صوتی ۶۰ بایتی (پس از آنکه \acro{ROHC} سرآیند ۶۰ بایتی \lr{IP/UDP/RTP} را به چند بایت فشرده کند، $L$ باز هم
کوچک‌تر است): $\eta = 85.7\,\%$. حفاظت یکپارچگی یک \lr{MAC-I} چهاربایتی به \acro{PDCP} می‌افزاید.

\subsection{نرخ فرانت‌هال}\label{sec:fh}
\paragraph{مدل} با گزینهٔ تقسیم ۸، فرانت‌هال نمونه‌های \acro{IQ} حوزهٔ زمان را حمل می‌کند: نرخ نمونه‌برداری $f_s$، $b$ بیت برای هر مؤلفهٔ
\lr{I} و \lr{Q}، $N_{\mathrm{ant}}$ آنتن و کدگذاری خط $c$:
\begin{equation}\label{eq:cpri}
  R_{\mathrm{FH},8} = 2\, f_s\, b\, N_{\mathrm{ant}}\, c .
\end{equation}
با تقسیم ۷٫۲، فرانت‌هال نمونه‌های حوزهٔ بسامد عناصر منبع اشغال‌شده را برای هر لایه حمل می‌کند، تقریباً
$R_{\mathrm{FH},7.2} \approx 12\,N_{\mathrm{PRB}}\,N_{\mathrm{sym}}\,N_{\mathrm{layers}}\,2b'$ که در آن $N_{\mathrm{sym}}$ تعداد نماد در ثانیه و $b'$
پهنای بیت فشرده‌شده است.
\paragraph{مثال} \acro{NR} با ۱۰۰ مگاهرتز در ۳۰ کیلوهرتز: $f_s = 122.88$ مگانمونه بر ثانیه، $b = 16$، ۴ آنتن، $c = 66/64$:
$R_{\mathrm{FH},8} = 16.2$ گیگابیت بر ثانیه؛ با ۶۴ آنتن ۲۶۰ گیگابیت بر ثانیه -- غیرعملی. تقسیم ۷٫۲ با ۲۷۳ \acro{PRB}، ۲۸\,۰۰۰ نماد بر ثانیه، ۴ لایه و
$b' = 9$: حدود ۶٫۶ گیگابیت بر ثانیه، که با تعداد لایه‌ها و نه آنتن‌ها مقیاس می‌یابد. تقسیم لایهٔ بالاتر (گزینهٔ ۲) تقریباً نرخ دادهٔ کاربر را
حمل می‌کند (حدود ۲٫۵ گیگابیت بر ثانیه برای حامل درس ۱۴).
\paragraph{تفسیر} \acro{MIMO} انبوه تنها با تقسیم‌هایی عملی است که شکل‌دهی پرتو را در \acro{RU} انجام می‌دهند؛ انتخاب تقسیم، معاوضه‌ای میان
تمرکز و هزینهٔ انتقال است.

\subsection{بودجهٔ تأخیر}\label{sec:budget}
\paragraph{مدل} تأخیر یک‌طرفهٔ صفحهٔ کاربر میان \acro{UE} و یک کارساز کاربرد مجموع اجزای آن است:
\begin{equation}\label{eq:budget}
  T = T_{\mathrm{radio}} + T_{\mathrm{gNB}} + T_{\mathrm{FH}} + T_{\mathrm{BH}} + T_{\mathrm{UPF}} + T_{\mathrm{N6}} .
\end{equation}
\paragraph{مثال} با مقادیر نمونه (\cref{fig:budget}): رادیو ۲ میلی‌ثانیه، پردازش \acro{gNB} ۰٫۵ میلی‌ثانیه، فرانت‌هال ۰٫۱ میلی‌ثانیه، \acro{UPF} ۰٫۱
میلی‌ثانیه؛ یک \acro{UPF} لبه با ۰٫۲ میلی‌ثانیه انتقال و ۰٫۱ میلی‌ثانیه تا کارساز ۳٫۰ میلی‌ثانیه می‌دهد، و یک \acro{UPF} مرکزی در فاصلهٔ ۶۰۰
کیلومتر (۳ میلی‌ثانیه انتقال) با ۱ میلی‌ثانیه تا کارساز ۶٫۷ میلی‌ثانیه. تنها جانمایی لبه هدف یک‌طرفهٔ ۴ میلی‌ثانیه را برآورده می‌کند.

\begin{acnfigure}
  \acnfig{lecture-23/latency-budget}
  \caption{بودجهٔ تأخیر یک‌طرفهٔ صفحهٔ کاربر با \acro{UPF} لبه و مرکزی (مقادیر نمونه، \cref{eq:budget}).}
  \label{fig:budget}
\end{acnfigure}

\subsection{قابلیت اطمینان با تکرار}\label{sec:dup}
\paragraph{مدل} اگر هر یک از دو مسیر مستقل یک بسته را با احتمال $p_1$ و $p_2$ از دست بدهد، تکرار تنها وقتی بسته را از دست می‌دهد که هر دو
شکست بخورند:
\begin{equation}\label{eq:dup}
  P_{\mathrm{loss}} = p_1\,p_2 .
\end{equation}
\paragraph{مثال} $p_1 = p_2 = 10^{-3}$ مقدار $10^{-6}$ را می‌دهد و یک هدف \acro{URLLC} را برآورده می‌کند که یک مسیر تنها فقط با یک \acro{MCS} بسیار
مقاوم‌تر (و پرهزینه‌تر) یا ارسال‌های مجدد بیشتر به آن می‌رسید -- که ممکن است در بودجهٔ تأخیر نگنجد.
\paragraph{محدودیت} اتلاف‌های همبسته (هر دو حامل از یک سایت که با یک مانع مسدود می‌شوند) فرض استقلال را نقض می‌کنند؛ اتصال دوگانه با گره‌های
جدا از هم از نظر جغرافیایی از تکرار در سطح حامل مقاوم‌تر است.

\subsection{پنجرهٔ شمارهٔ ترتیب}\label{sec:sn}
\paragraph{مدل} با یک \acro{SN} $s$ بیتی، حداکثر $2^{s-1}$ \acro{PDU} می‌توانند در جریان باشند. با نرخ بستهٔ $P$ پنجره این مدت را پوشش می‌دهد:
\begin{equation}\label{eq:window}
  T_{\mathrm{win}} = \frac{2^{s-1}}{P} .
\end{equation}
\paragraph{مثال} ۲ گیگابیت بر ثانیه با بسته‌های ۱۵۰۰ بایتی: $P = 1.67\times10^5$ بسته بر ثانیه. یک \acro{SN} ۱۲ بیتی ۱۲٫۳ میلی‌ثانیه را پوشش می‌دهد
-- کوتاه‌تر از تأخیرهای \acro{HARQ} و \acro{ARQ} به‌علاوهٔ بازترتیب‌دهی؛ یک \acro{SN} ۱۸ بیتی ۰٫۷۹ ثانیه. بنابراین حامل‌های پرگذردهی از \acro{SN}های ۱۸
بیتی استفاده می‌کنند.

% ---------------------------------------------------------------------
\section{معماری شبکه و نمودارها}
یک استقرار نوعی نسل پنجم \acro{RU}ها را در سایت‌های آنتن، \acro{DU}ها را در سایت‌های سلول یا مراکز نزدیک (متصل با فرانت‌هال ۷٫۲)، \lr{CU-UP}ها را
همراه با \acro{UPF}های لبه در مراکز دادهٔ منطقه‌ای، و \acro{UPF}های مرکزی را برای ترافیک اینترنت ترکیب می‌کند. هر عنصر را می‌توان روی زیرساخت
مشترک مجازی کرد (\acro{vRAN}، درس ۱۶) و \acro{SMF} برای هر نشست \acro{UPF}ها را زنجیر می‌کند (\cref{fig:stack,fig:splits,fig:upf}).

% ---------------------------------------------------------------------
\section{رویه‌های گام‌به‌گام}

\begin{procedurebox}{دنبال‌کردن یک بستهٔ فروسو}\label{proc:dl}
\begin{steps}
  \item بسته روی \lr{N6} به \acro{UPF} از نوع \acro{PSA} می‌رسد؛ یک \acro{PDR} با آن منطبق می‌شود (نشانی \acro{UE} و فیلتر \acro{SDF})، \acro{QER} نرخ‌ها را
        اعمال می‌کند، \acro{URR} آن را می‌شمارد و \acro{FAR} آن را با \acro{QFI} در \acro{GTP-U} به سوی گام بعد (\lr{N9} یا \lr{N3}) کپسوله می‌کند.
  \item \lr{CU-UP} آن را روی \lr{N3} دریافت می‌کند، \acro{SDAP} \acro{QFI} را به یک \acro{DRB} نگاشت می‌کند، \acro{PDCP} سرآیندها را فشرده، رمز و
        شماره‌گذاری می‌کند و آن را در \acro{GTP-U} روی \lr{F1-U} می‌فرستد.
  \item \acro{RLC} در \acro{DU} آن را طبق زمان‌بندی \acro{MAC} قطعه‌بندی می‌کند؛ \acro{MAC} آن را در یک بلوک انتقال تسهیم می‌کند؛ \acro{PHY} آن را
        با \acro{HARQ} ارسال می‌کند.
  \item لایه‌های \acro{UE} فرایند را وارونه می‌کنند؛ \acro{PDCP} بازترتیب می‌دهد و به‌ترتیب به \acro{SDAP} و لایهٔ \acro{IP} تحویل می‌دهد.
\end{steps}
\end{procedurebox}

\begin{procedurebox}{انتخاب تقسیم \lr{RAN}}\label{proc:split}
\begin{steps}
  \item پهنای باند، تعداد آنتن‌ها و لایه‌ها در هر سلول را تعیین و نرخ فرانت‌هال را برای تقسیم‌های نامزد محاسبه کنید (\cref{eq:cpri}).
  \item با انتقال موجود (فیبر، ریزموج) و تأخیر آن مقایسه کنید؛ تقسیم‌های پایین‌تر به تأخیر فرانت‌هال زیر یک میلی‌ثانیه نیاز دارند.
  \item مزایای تمرکز (هماهنگی، تجمیع، سایت‌های ساده‌تر) را در برابر هزینهٔ انتقال بسنجید.
  \item تقسیم را بر حسب نوع سایت انتخاب کنید (مثلاً ۷٫۲ برای \acro{MIMO} انبوه شهری، گزینهٔ ۲ برای سایت‌های دورافتاده با انتقال محدود).
\end{steps}
\end{procedurebox}

% ---------------------------------------------------------------------
\section{مثال‌های حل‌شده}

\begin{examplebox}{سربار برای بسته‌های کوچک}\label{ex:oh}
\textbf{مسئله.} حسگرهای اینترنت اشیا محموله‌های ۲۰ بایتی را در \lr{IPv6/UDP} بدون فشرده‌سازی سرآیند می‌فرستند. بازدهی روی رادیو را محاسبه کنید،
و نیز وقتی \acro{ROHC} سرآیند \lr{IPv6/UDP} را به ۴ بایت فشرده می‌کند.

\textbf{حل.} بدون \acro{ROHC}: $L = 20 + 48 = 68$، بازدهی داده $20/(68+10) = 25.6\,\%$. با \acro{ROHC}: $L = 24$، $20/34 = 58.8\,\%$. فشرده‌سازی سرآیند
بیشترین اهمیت را برای بسته‌های کوچک دارد.
\end{examplebox}

\begin{examplebox}{فرانت‌هال برای \lr{MIMO} انبوه}\label{ex:fh}
\textbf{مسئله.} یک \acro{RU} با ۶۴ آنتن و ۱۰۰ مگاهرتز به ۱۶ لایه خدمت می‌دهد. فرانت‌هال گزینهٔ ۸ و گزینهٔ ۷٫۲ را مقایسه کنید (فرض‌های \cref{sec:fh}).

\textbf{حل.} گزینهٔ ۸: ۲۶۰ گیگابیت بر ثانیه. گزینهٔ ۷٫۲: $273\times12\times28\,000\times16\times18 = 26.4$ گیگابیت بر ثانیه -- ده برابر کمتر، و
روی یک پیوند اترنت ۲۵ یا ۵۰ گیگابیتی شدنی.
\end{examplebox}

\begin{examplebox}{بودجهٔ تکرار}\label{ex:dup}
\textbf{مسئله.} هر مسیر \acro{BLER} برابر $10^{-2}$ در هر ارسال دارد و می‌تواند در بودجه ۲ ارسال انجام دهد. یک مسیر تنها به چه اتلاف باقی‌مانده‌ای
می‌رسد و تکرار به چه مقداری؟

\textbf{حل.} مسیر تنها: $(10^{-2})^2 = 10^{-4}$. با تکرار روی دو مسیر مستقل: $10^{-4}\times10^{-4} = 10^{-8}$ (\cref{eq:dup}) -- فراتر از اهداف
$10^{-5}$ در \acro{URLLC}، با حاشیه برای همبستگی.
\end{examplebox}

% ---------------------------------------------------------------------
\section{مقایسهٔ فناوری‌ها}
\begin{table}[htbp]
\centering\small
\caption{صفحهٔ کاربر در \lr{LTE/EPC} و در \lr{NR/5GC}.}\label{tab:cmp}
\rowcolors{2}{acnmist}{white}
\begin{tabularx}{\linewidth}{@{}l Y Y@{}}
\toprule
\rowcolor{white} & \lr{LTE / EPC} & \lr{NR / 5GC} \\
\midrule
لایه‌های رادیویی   & \acro{PDCP}، \acro{RLC} (با بازترتیب‌دهی)، \acro{MAC} & \acro{SDAP}، \acro{PDCP}، \acro{RLC} (بدون بازترتیب‌دهی)، \acro{MAC} \\
تقسیم \acro{RAN}    & \acro{eNB} یکپارچه (\acro{CPRI} تا \acro{RRH}) & \acro{CU}/\acro{DU} (\lr{F1})، \lr{CU-CP}/\lr{CU-UP} (\lr{E1})، تقسیم لایهٔ پایین ۷٫۲ \\
صفحهٔ کاربر هسته    & \lr{S-GW} + \lr{P-GW} (\acro{CUPS} اختیاری) & زنجیرهٔ \acro{UPF}ها (\lr{I-UPF}، \lr{UL CL}، \acro{PSA}) با کنترل \acro{PFCP} \\
\acro{QoS} روی \lr{N3/S1-U} & تونل برای هر حامل & تونل برای هر نشست، \acro{QFI} برای هر بسته \\
قابلیت اطمینان      & \acro{HARQ}، \acro{ARQ} & + تکرار \acro{PDCP}، چنداتصالی \\
طول \acro{SN}       & \acro{PDCP} تا ۱۸ بیت (بعدها) & \acro{PDCP} ۱۲/۱۸، \lr{RLC AM} ۱۲/۱۸ بیت \\
\bottomrule
\end{tabularx}
\end{table}

\paragraph{برنامه‌پذیر از سر تا ته} از \acro{UPF}های برنامه‌ریزی‌شده با \lr{P4} تا \acro{CU}/\acro{DU} مجازی، صفحهٔ کاربر نسل پنجم سرتاسر نرم‌افزارمحور
است -- و این کارایی آن را به موضوعی از بودجه‌های مهندسی (چرخه به ازای هر بسته، میکروثانیه به ازای هر گام) تبدیل می‌کند و نه جعبه‌های ثابت.

% ---------------------------------------------------------------------
\section{نکات کلیدی}
\begin{takeawaybox}
\begin{itemize}
  \item دادهٔ کاربر: \acro{UE} (\lr{SDAP/PDCP/RLC/MAC/PHY}) $\to$ \acro{DU} (\lr{RLC/MAC/PHY}) $\to$ \lr{F1-U} $\to$ \lr{CU-UP} (\lr{PDCP/SDAP}) $\to$ \lr{N3}
        $\to$ \acro{UPF}(ها) $\to$ \lr{N6}.
  \item تقسیم لایهٔ بالاتر (گزینهٔ ۲) حدود نرخ کاربر نیاز دارد؛ گزینهٔ ۸ با تعداد آنتن‌ها مقیاس می‌یابد ($2f_s b N_{\mathrm{ant}} c$)؛ ۷٫۲ با تعداد لایه‌ها.
  \item پردازش \acro{UPF}: تطابق \acro{PDR} $\to$ \acro{QER} $\to$ \acro{URR} $\to$ \acro{FAR}، که \acro{SMF} روی \lr{N4} برنامه‌ریزی می‌کند.
  \item تأخیر $= \sum$ اجزا؛ \acro{UPF}های لبه چند میلی‌ثانیه انتقال را حذف می‌کنند.
  \item تکرار \acro{PDCP} مقدار $P_{\mathrm{loss}} = p_1p_2$ را می‌دهد؛ \acro{SN}های بلند پنجرهٔ $2^{s-1}/P$ را بالاتر از تأخیرهای بازترتیب‌دهی نگه می‌دارند.
\end{itemize}
\end{takeawaybox}

% ---------------------------------------------------------------------
\section{تمرین‌ها}

\begin{exercisebox}{1}{کدام لایه؟}\label{exr:layer}
کدام لایه این کارها را انجام می‌دهد: (الف) رمزنگاری؛ (ب) قطعه‌بندی؛ (ج) نگاشت جریان \acro{QoS} به \acro{DRB}؛ (د) \acro{HARQ}؛ (هـ) حذف تکراری‌ها؛
(و) \acro{ARQ}؟
\end{exercisebox}

\begin{exercisebox}{1}{چرا تقسیم؟}\label{exr:why}
دو مزیت و دو عیب متمرکزکردن \acro{CU} برای \acro{DU}های بسیار را بیان کنید.
\end{exercisebox}

\begin{exercisebox}{2}{فرانت‌هال}\label{exr:fh}
یک حامل ۴۰ مگاهرتزی در ۳۰ کیلوهرتز با ۶۱٫۴۴ مگانمونه بر ثانیه نمونه‌برداری می‌شود. فرانت‌هال گزینهٔ ۸ را برای ۸ آنتن با نمونه‌های ۱۶ بیتی و
کدگذاری ۶۶/۶۴ محاسبه کنید. آیا روی یک پیوند ۲۵ گیگابیتی اجرا می‌شود؟
\end{exercisebox}

\begin{exercisebox}{2}{پنجرهٔ \lr{SN}}\label{exr:sn}
یک حامل ۵ گیگابیت بر ثانیه بستهٔ ۱۲۰۰ بایتی حمل می‌کند. پنجرهٔ پوشش‌داده‌شده با \acro{SN}های ۱۲ و ۱۸ بیتی \acro{PDCP} را محاسبه کنید. اگر زمان‌سنج
بازترتیب‌دهی بتواند به ۱۰۰ میلی‌ثانیه برسد، کدام لازم است؟
\end{exercisebox}

\begin{exercisebox}{3}{جانمایی \lr{UPF}ها}\label{exr:place}
یک اپراتور یک مرکز دادهٔ مرکزی (۸۰۰ کیلومتر از دورترین سلول‌ها) دارد و می‌تواند سایت‌های منطقه‌ای حداکثر در ۱۵۰ کیلومتری هر سلول بسازد. یک خدمت به
تأخیر یک‌طرفهٔ ۵ میلی‌ثانیه نیاز دارد، با رادیو ۲٫۵ میلی‌ثانیه، \acro{gNB} ۰٫۵ میلی‌ثانیه و \acro{UPF} و کارساز ۰٫۴ میلی‌ثانیه. با فرض ۵ میکروثانیه بر
کیلومتر برای انتقال، کدام جانمایی کار می‌کند؟ هزینه را بحث کنید.
\end{exercisebox}

\begin{exercisebox}{3}{تکرار همبسته}\label{exr:corr}
دو مسیر هر کدام بسته‌ها را با احتمال $10^{-3}$ از دست می‌دهند، اما با احتمال $q$ با هم مسدود می‌شوند (علت مشترک) و در این حالت هر دو از دست
می‌دهند. اتلاف را به‌صورت $P = q + (1-q)p^2$ با $p = 10^{-3}$ مدل کنید و بزرگ‌ترین $q$ را بیابید که $P \le 10^{-5}$. این برای استقرار چه معنایی دارد؟
\end{exercisebox}

\subsection*{پاسخ‌ها}
\addcontentsline{toc}{subsection}{پاسخ‌ها}

\begin{solutionbox}{\ref{exr:layer}}
(الف) \acro{PDCP}؛ (ب) \acro{RLC}؛ (ج) \acro{SDAP}؛ (د) \acro{MAC}؛ (هـ) \acro{PDCP}؛ (و) \acro{RLC} (حالت تصدیق‌شده).
\end{solutionbox}

\begin{solutionbox}{\ref{exr:why}}
مزایا: تجمیع پردازش \acro{PDCP}/\acro{RRC} (تسهیم آماری)، هماهنگی و جابه‌جایی آسان‌تر میان سلول‌های یک \acro{CU}، سایت‌های سلولی ساده‌تر و ارزان‌تر،
ارتقای نرم‌افزار در یک جا. معایب: نیاز به انتقال میان \acro{CU} و \acro{DU} (پهنای باند و تأخیر)، دامنهٔ خرابی بزرگ‌تر (قطعی یک \acro{CU} سلول‌های بسیاری
را متأثر می‌کند)، و کارکردهای حساس به تأخیر را نمی‌توان فراتر از یک فاصله متمرکز کرد.
\end{solutionbox}

\begin{solutionbox}{\ref{exr:fh}}
$2\times61.44\times10^6\times16\times8\times66/64 = 16.2$ گیگابیت بر ثانیه: برای یک حامل در پیوند ۲۵ گیگابیتی جا می‌شود، اما نه برای دو حامل
این‌چنین.
\end{solutionbox}

\begin{solutionbox}{\ref{exr:sn}}
$P = 5\times10^9/9600 = 5.2\times10^5$ بسته بر ثانیه. ۱۲ بیتی: $2048/P = 3.9$ میلی‌ثانیه؛ ۱۸ بیتی: $131\,072/P = 0.25$ ثانیه. با زمان‌سنج بازترتیب‌دهی
۱۰۰ میلی‌ثانیه، تنها \acro{SN} ۱۸ بیتی امن است.
\end{solutionbox}

\begin{solutionbox}{\ref{exr:place}}
بخش ثابت $2.5 + 0.5 + 0.4 = 3.4$ میلی‌ثانیه است و ۱٫۶ میلی‌ثانیه برای انتقال باقی می‌ماند: $D \le 1.6\ \text{ms}/5\ \mu\text{s/km} = 320$ کیلومتر.
سایت مرکزی (۸۰۰ کیلومتر، ۴ میلی‌ثانیه) شکست می‌خورد؛ سایت‌های منطقه‌ای (۱۵۰ کیلومتر، ۰٫۷۵ میلی‌ثانیه) کار می‌کنند. هزینه: سایت‌های \acro{UPF} بیشتر
(سخت‌افزار، عملیات)، همتاسازی یا نهان‌سازی منطقه‌ای برای خدمت‌رسانی محلی محتوا -- که باید با درآمد خدمات حساس به تأخیر توجیه شود.
\end{solutionbox}

\begin{solutionbox}{\ref{exr:corr}}
$q + (1-q)10^{-6} \le 10^{-5} \Rightarrow q \le 9\times10^{-6}$ (تقریباً). مسدودشدن با علت مشترک باید بسیار نادر باشد: مسیرها باید از سایت‌ها و
جهت‌های متفاوت بیایند (اتصال دوگانه)، نه صرفاً از حامل‌های متفاوت یک آنتن.
\end{solutionbox}

% ---------------------------------------------------------------------
\section{مطالعهٔ بیشتر}
معماری صفحهٔ کاربر \acro{NR} در \lr{TS 38.300} \cite{ts38300} و \acro{RAN} تقسیم‌شده در \lr{TS 38.401} \cite{ts38401} شرح داده شده‌اند؛ \acro{PFCP} در
\lr{TS 29.244} \cite{ts29244}؛ \acro{GTP-U} در \lr{TS 29.281} \cite{ts29281}؛ کارکردها و توپولوژی‌های \acro{UPF} در \lr{TS 23.501} \cite{ts23501}. دالمن و
همکاران \cite{dahlman2020nr} پروتکل‌های لایهٔ ۲ در \acro{NR} را توضیح می‌دهند؛ بوسهارت و همکاران \cite{bosshart2014p4} زبان \lr{P4} را که در
\acro{UPF}های برنامه‌پذیر به کار می‌رود معرفی می‌کنند.

\begin{latin}
\printbibliography[heading=subbibliography,title={References}]
\end{latin}
\end{document}
